\documentclass[11pt, letterpaper]{article}
\usepackage[utf8]{inputenc}
\usepackage{amsmath, amssymb, amsthm}
\usepackage{geometry}
\usepackage{graphicx}
\usepackage{fancyhdr}
\usepackage{enumitem}
\usepackage{tikz}
\usepackage{xcolor}

% Page Geometry
\geometry{margin=1in}

% Header/Footer
\pagestyle{fancy}
\fancyhf{}
\lhead{Homework 14.5: The Chain Rule}
\rhead{Tashfeen Omran}
\rfoot{Page \thepage}

% Custom Commands for Formatting
\newcommand{\sectionbreak}{\clearpage}
\newtheorem{theorem}{Theorem}
\theoremstyle{definition}
\newtheorem{definition}{Definition}
\newtheorem{example}{Example}

\title{\textbf{Comprehensive Study Guide \& Solutions}\\ Calculus III: The Multivariable Chain Rule}
\author{Tashfeen Omran}
\date{Spring 2026}

\begin{document}

\maketitle
\tableofcontents
\newpage

% =======================================================
% PART 1: INTRODUCTION
% =======================================================
\section{Part 1: Introduction, Context, and Prerequisites}

\subsection{Core Concepts}
The \textbf{Multivariable Chain Rule} is the primary tool for calculating how a function of several variables changes when those input variables are themselves changing. 

In single-variable calculus, if $y = f(x)$ and $x = g(t)$, the Chain Rule tells us that $\frac{dy}{dt} = \frac{dy}{dx} \cdot \frac{dx}{dt}$. In multivariable calculus, the concept is similar but involves summing the contributions from every variable.

If $z = f(x, y)$ is a function where $x$ and $y$ both depend on a parameter $t$, a change in $t$ causes a change in $x$ (affecting $z$) \textit{and} a change in $y$ (also affecting $z$). The total rate of change is the sum of these effects.

\subsection{Intuition and Tree Diagrams}
The most intuitive way to visualize this is using a \textbf{Tree Diagram}.
\begin{itemize}
    \item The top of the tree is the dependent variable (e.g., $z$).
    \item The middle branches are the intermediate variables (e.g., $x, y$).
    \item The bottom roots are the independent variables (e.g., $t$, or $u, v$).
\end{itemize}
To find a derivative, you trace every path from the top variable to the bottom variable, multiplying the derivatives along each branch, and then adding the results of all paths together.

\subsection{Historical Context}
The Chain Rule in multiple dimensions was formalized as calculus expanded from studying curves (1D) to surfaces and physics in 3D space (18th-19th century). Mathematicians like Euler and Lagrange required these tools to describe "variational calculus" and fluid dynamics, where a particle's position $(x, y, z)$ changes over time $t$, but the field it moves through (like temperature or potential energy) varies with position.

\subsection{Key Formulas}

\textbf{Case 1: One Independent Variable} ($z$ depends on $x,y$; $x,y$ depend on $t$)
\[ \frac{dz}{dt} = \frac{\partial z}{\partial x}\frac{dx}{dt} + \frac{\partial z}{\partial y}\frac{dy}{dt} \]

\textbf{Case 2: Several Independent Variables} ($z$ depends on $x,y$; $x,y$ depend on $s,t$)
\[ \frac{\partial z}{\partial s} = \frac{\partial z}{\partial x}\frac{\partial x}{\partial s} + \frac{\partial z}{\partial y}\frac{\partial y}{\partial s} \]
\[ \frac{\partial z}{\partial t} = \frac{\partial z}{\partial x}\frac{\partial x}{\partial t} + \frac{\partial z}{\partial y}\frac{\partial y}{\partial t} \]

\textbf{Implicit Differentiation Theorem}
If $F(x, y) = 0$ defines $y$ implicitly as a function of $x$:
\[ \frac{dy}{dx} = -\frac{F_x}{F_y} = -\frac{\partial F / \partial x}{\partial F / \partial y} \]

\subsection{Prerequisites}
\begin{itemize}
    \item \textbf{Partial Differentiation:} Ability to hold one variable constant while differentiating another.
    \item \textbf{Single-Variable Chain Rule:} $(f(g(x)))' = f'(g(x))g'(x)$.
    \item \textbf{Product and Quotient Rules:} Essential for complex algebraic manipulations.
    \item \textbf{Trigonometric Derivatives:} $\frac{d}{dx}\tan(x) = \sec^2(x)$, $\frac{d}{dx}\arctan(x) = \frac{1}{1+x^2}$, etc.
\end{itemize}

\newpage

% =======================================================
% PART 2: SOLUTIONS
% =======================================================
\section{Part 2: Detailed Homework Solutions}

\subsection*{Problem 1: Chain Rule vs. Substitution}
\textbf{Given:} $z = x^2y + xy^2$, $x = 9t$, $y = t^2$.

\subsubsection*{Part A: Use the Chain Rule}
Formula: $\frac{dz}{dt} = \frac{\partial z}{\partial x}\frac{dx}{dt} + \frac{\partial z}{\partial y}\frac{dy}{dt}$.

Step 1: Find Partial Derivatives of $z$.
\[ \frac{\partial z}{\partial x} = 2xy + y^2 \]
\[ \frac{\partial z}{\partial y} = x^2 + 2xy \]

Step 2: Find Derivatives of $x$ and $y$ with respect to $t$.
\[ \frac{dx}{dt} = 9, \quad \frac{dy}{dt} = 2t \]

Step 3: Apply the Chain Rule.
\[ \frac{dz}{dt} = (2xy + y^2)(9) + (x^2 + 2xy)(2t) \]

Step 4: Substitute expressions for $x$ and $y$ in terms of $t$ to finalize the answer.
Substitute $x=9t$ and $y=t^2$:
\begin{align*} 
\frac{dz}{dt} &= [2(9t)(t^2) + (t^2)^2](9) + [(9t)^2 + 2(9t)(t^2)](2t) \\
&= [18t^3 + t^4](9) + [81t^2 + 18t^3](2t) \\
&= 162t^3 + 9t^4 + 162t^3 + 36t^4 \\
&= 45t^4 + 324t^3 
\end{align*}

\subsubsection*{Part B: Substitution Method}
Step 1: Write $z$ as a function of $t$.
\[ z = (9t)^2(t^2) + (9t)(t^2)^2 = 81t^2 \cdot t^2 + 9t \cdot t^4 = 81t^4 + 9t^5 \]

Step 2: Differentiate with respect to $t$.
\[ \frac{dz}{dt} = \frac{d}{dt}(81t^4 + 9t^5) = 324t^3 + 45t^4 \]

\textbf{Agreement Check:} Both methods yield $45t^4 + 324t^3$.

\textbf{Final Answer:} $\boxed{45t^4 + 324t^3}$

% -------------------------------------------------------

\subsection*{Problem 2: Multiple Partial Derivatives}
\textbf{Given:} $z = xy^5 - x^2y$, $x = t^2 + 1$, $y = t^2 - 1$.

\textbf{Find:} $\frac{\partial z}{\partial x}, \frac{\partial z}{\partial y}, \frac{dx}{dt}, \frac{dy}{dt}, \frac{dz}{dt}$.

1. $\frac{\partial z}{\partial x}$:
\[ \frac{\partial}{\partial x}(xy^5 - x^2y) = y^5 - 2xy \]

2. $\frac{\partial z}{\partial y}$:
\[ \frac{\partial}{\partial y}(xy^5 - x^2y) = 5xy^4 - x^2 \]

3. $\frac{dx}{dt}$:
\[ \frac{d}{dt}(t^2+1) = 2t \]

4. $\frac{dy}{dt}$:
\[ \frac{d}{dt}(t^2-1) = 2t \]

5. Use the Chain Rule for $\frac{dz}{dt}$:
\[ \frac{dz}{dt} = \frac{\partial z}{\partial x}\frac{dx}{dt} + \frac{\partial z}{\partial y}\frac{dy}{dt} \]
Substitute the parts found above:
\[ \frac{dz}{dt} = (y^5 - 2xy)(2t) + (5xy^4 - x^2)(2t) \]
(Note: The problem usually accepts the answer in mixed variables unless specified "in terms of $t$". Based on standard WebAssign formats, you may need to substitute $t$. Let's provide the substituted form for completeness).

Substitution:
$y^5 - 2xy = (t^2-1)^5 - 2(t^2+1)(t^2-1) = (t^2-1)^5 - 2(t^4-1)$.
$5xy^4 - x^2 = 5(t^2+1)(t^2-1)^4 - (t^2+1)^2$.
\[ \frac{dz}{dt} = 2t \left[ (t^2-1)^5 - 2(t^4-1) + 5(t^2+1)(t^2-1)^4 - (t^2+1)^2 \right] \]

\textbf{Final Answer:}
$\frac{\partial z}{\partial x} = \boxed{y^5 - 2xy}$ \\
$\frac{\partial z}{\partial y} = \boxed{5xy^4 - x^2}$ \\
$\frac{dx}{dt} = \boxed{2t}$ \\
$\frac{dy}{dt} = \boxed{2t}$ \\
$\frac{dz}{dt} = \boxed{2t(y^5 - 2xy + 5xy^4 - x^2)}$ (or fully substituted form).

% -------------------------------------------------------

\subsection*{Problem 3: Trigonometric Chain Rule}
\textbf{Given:} $z = \sin(x)\cos(y)$, $x = \sqrt{t} = t^{1/2}$, $y = \frac{5}{t} = 5t^{-1}$.

\textbf{Find:} $\frac{dz}{dt}$.

Step 1: Partials of $z$.
\[ \frac{\partial z}{\partial x} = \cos(x)\cos(y) \]
\[ \frac{\partial z}{\partial y} = -\sin(x)\sin(y) \]

Step 2: Derivatives of $x, y$.
\[ \frac{dx}{dt} = \frac{1}{2}t^{-1/2} = \frac{1}{2\sqrt{t}} \]
\[ \frac{dy}{dt} = -5t^{-2} = -\frac{5}{t^2} \]

Step 3: Assemble.
\[ \frac{dz}{dt} = (\cos x \cos y)\left(\frac{1}{2\sqrt{t}}\right) + (-\sin x \sin y)\left(-\frac{5}{t^2}\right) \]
\[ \frac{dz}{dt} = \frac{\cos(\sqrt{t})\cos(5/t)}{2\sqrt{t}} + \frac{5\sin(\sqrt{t})\sin(5/t)}{t^2} \]

\textbf{Final Answer:} $\boxed{\frac{\cos(\sqrt{t})\cos(5/t)}{2\sqrt{t}} + \frac{5\sin(\sqrt{t})\sin(5/t)}{t^2}}$

% -------------------------------------------------------

\subsection*{Problem 4: Inverse Trig and Roots}
\textbf{Given:} $z = \sqrt{7 + xy}$, $x = \tan(t)$, $y = \arctan(t)$.

\textbf{Find:} $\frac{dz}{dt}$.

Step 1: Partials of $z = (7+xy)^{1/2}$.
\[ \frac{\partial z}{\partial x} = \frac{1}{2}(7+xy)^{-1/2} \cdot y = \frac{y}{2\sqrt{7+xy}} \]
\[ \frac{\partial z}{\partial y} = \frac{1}{2}(7+xy)^{-1/2} \cdot x = \frac{x}{2\sqrt{7+xy}} \]

Step 2: Derivatives of $x, y$.
\[ \frac{dx}{dt} = \sec^2(t) \]
\[ \frac{dy}{dt} = \frac{1}{1+t^2} \]

Step 3: Chain Rule.
\[ \frac{dz}{dt} = \frac{y}{2\sqrt{7+xy}}\sec^2(t) + \frac{x}{2\sqrt{7+xy}}\left(\frac{1}{1+t^2}\right) \]
Factor out the common term:
\[ \frac{dz}{dt} = \frac{1}{2\sqrt{7+xy}} \left( y\sec^2(t) + \frac{x}{1+t^2} \right) \]

Substitute $x, y$ back in:
\[ \frac{dz}{dt} = \frac{1}{2\sqrt{7+\tan(t)\arctan(t)}} \left( \arctan(t)\sec^2(t) + \frac{\tan(t)}{1+t^2} \right) \]

\textbf{Final Answer:} $\boxed{\frac{\arctan(t)\sec^2(t) + \frac{\tan(t)}{1+t^2}}{2\sqrt{7+\tan(t)\arctan(t)}}}$

% -------------------------------------------------------

\subsection*{Problem 5: Two Independent Variables}
\textbf{Given:} $z = x^2 + y^2$, $x = 5s + 8t$, $y = s + t$.

\textbf{Find:} $\frac{\partial z}{\partial s}$ and $\frac{\partial z}{\partial t}$.

Partials of $z$: $\frac{\partial z}{\partial x} = 2x$, $\frac{\partial z}{\partial y} = 2y$.

Derivatives of $x$ and $y$:
$\frac{\partial x}{\partial s} = 5, \quad \frac{\partial x}{\partial t} = 8$.
$\frac{\partial y}{\partial s} = 1, \quad \frac{\partial y}{\partial t} = 1$.

\textbf{Calculate $\frac{\partial z}{\partial s}$:}
\[ \frac{\partial z}{\partial s} = \frac{\partial z}{\partial x}\frac{\partial x}{\partial s} + \frac{\partial z}{\partial y}\frac{\partial y}{\partial s} = (2x)(5) + (2y)(1) = 10x + 2y \]
Substitute $s, t$:
$10(5s+8t) + 2(s+t) = 50s + 80t + 2s + 2t = 52s + 82t$.

\textbf{Calculate $\frac{\partial z}{\partial t}$:}
\[ \frac{\partial z}{\partial t} = \frac{\partial z}{\partial x}\frac{\partial x}{\partial t} + \frac{\partial z}{\partial y}\frac{\partial y}{\partial t} = (2x)(8) + (2y)(1) = 16x + 2y \]
Substitute $s, t$:
$16(5s+8t) + 2(s+t) = 80s + 128t + 2s + 2t = 82s + 130t$.

\textbf{Alternative Method (Substitution):}
$z = (5s+8t)^2 + (s+t)^2 = (25s^2 + 80st + 64t^2) + (s^2 + 2st + t^2) = 26s^2 + 82st + 65t^2$.
$\frac{\partial z}{\partial s} = 52s + 82t$.
$\frac{\partial z}{\partial t} = 82s + 130t$.

\textbf{Do answers agree?} Yes.

\textbf{Final Answers:}
$\frac{\partial z}{\partial s} = \boxed{52s + 82t}$
$\frac{\partial z}{\partial t} = \boxed{82s + 130t}$

% -------------------------------------------------------

\subsection*{Problem 6: Power Rule Chain}
\textbf{Given:} $z = (x-y)^9$, $x = s^2t$, $y = st^2$.

Partials of $z$:
$\frac{\partial z}{\partial x} = 9(x-y)^8$.
$\frac{\partial z}{\partial y} = 9(x-y)^8(-1) = -9(x-y)^8$.

Derivatives of $x, y$:
$x_s = 2st, \quad x_t = s^2$.
$y_s = t^2, \quad y_t = 2st$.

\textbf{Find $\frac{\partial z}{\partial s}$:}
\[ \frac{\partial z}{\partial s} = 9(x-y)^8(2st) + [-9(x-y)^8](t^2) = 9(x-y)^8(2st - t^2) \]

\textbf{Find $\frac{\partial z}{\partial t}$:}
\[ \frac{\partial z}{\partial t} = 9(x-y)^8(s^2) + [-9(x-y)^8](2st) = 9(x-y)^8(s^2 - 2st) \]

\textbf{Final Answers:}
$\frac{\partial z}{\partial s} = \boxed{9(s^2t - st^2)^8(2st - t^2)}$
$\frac{\partial z}{\partial t} = \boxed{9(s^2t - st^2)^8(s^2 - 2st)}$

% -------------------------------------------------------

\subsection*{Problem 7: Quotient and Trig}
\textbf{Given:} $z = \tan(u/v)$, $u = 4s + 3t$, $v = 3s - 4t$.

Partials of $z$:
Let $A = u/v$. $\frac{\partial z}{\partial A} = \sec^2(u/v)$.
$\frac{\partial z}{\partial u} = \sec^2(u/v) \cdot \frac{1}{v}$.
$\frac{\partial z}{\partial v} = \sec^2(u/v) \cdot (-u v^{-2}) = -\frac{u}{v^2}\sec^2(u/v)$.

Derivatives of $u, v$:
$u_s = 4, \quad u_t = 3$.
$v_s = 3, \quad v_t = -4$.

\textbf{Find $\frac{\partial z}{\partial s}$:}
\begin{align*}
\frac{\partial z}{\partial s} &= \left(\frac{1}{v}\sec^2\frac{u}{v}\right)(4) + \left(-\frac{u}{v^2}\sec^2\frac{u}{v}\right)(3) \\
&= \sec^2\left(\frac{u}{v}\right) \left[ \frac{4}{v} - \frac{3u}{v^2} \right] = \sec^2\left(\frac{u}{v}\right) \left[ \frac{4v - 3u}{v^2} \right]
\end{align*}
Substitute $u, v$ into the bracket term:
$4v - 3u = 4(3s-4t) - 3(4s+3t) = 12s - 16t - 12s - 9t = -25t$.
Result: $\boxed{-\frac{25t}{(3s-4t)^2}\sec^2\left(\frac{4s+3t}{3s-4t}\right)}$

\textbf{Find $\frac{\partial z}{\partial t}$:}
\begin{align*}
\frac{\partial z}{\partial t} &= \left(\frac{1}{v}\sec^2\frac{u}{v}\right)(3) + \left(-\frac{u}{v^2}\sec^2\frac{u}{v}\right)(-4) \\
&= \sec^2\left(\frac{u}{v}\right) \left[ \frac{3}{v} + \frac{4u}{v^2} \right] = \sec^2\left(\frac{u}{v}\right) \left[ \frac{3v + 4u}{v^2} \right]
\end{align*}
Substitute $u, v$:
$3v + 4u = 3(3s-4t) + 4(4s+3t) = 9s - 12t + 16s + 12t = 25s$.
Result: $\boxed{\frac{25s}{(3s-4t)^2}\sec^2\left(\frac{4s+3t}{3s-4t}\right)}$

% -------------------------------------------------------

\subsection*{Problem 8: Abstract Chain Rule}
\textbf{Given:} $p(t) = f(g(t), h(t))$.
Data: $g(9)=6, g'(9)=-1, h(9)=7, h'(9)=8$.
$f_x(6,7)=5, f_y(6,7)=6$.

\textbf{Find:} $p'(9)$.

Formula: $p'(t) = f_x(g(t), h(t)) \cdot g'(t) + f_y(g(t), h(t)) \cdot h'(t)$.
At $t=9$:
$p'(9) = f_x(g(9), h(9)) g'(9) + f_y(g(9), h(9)) h'(9)$.
$p'(9) = f_x(6, 7)(-1) + f_y(6, 7)(8)$.
$p'(9) = (5)(-1) + (6)(8) = -5 + 48 = 43$.

\textbf{Final Answer:} $\boxed{43}$

% -------------------------------------------------------

\subsection*{Problem 9: Tree Diagram Identification}
\textbf{Given:} $u = f(x,y)$, $x=x(r,s,t)$, $y=y(r,s,t)$.
We need the chain rule for partials with respect to $r, s, t$.

For $\frac{\partial u}{\partial r}$:
Paths from $u$ to $r$ go through $x$ and $y$.
Path 1: $u \to x \to r \Rightarrow \frac{\partial u}{\partial x}\frac{\partial x}{\partial r}$.
Path 2: $u \to y \to r \Rightarrow \frac{\partial u}{\partial y}\frac{\partial y}{\partial r}$.
Sum: $\frac{\partial u}{\partial r} = \frac{\partial u}{\partial x}\frac{\partial x}{\partial r} + \frac{\partial u}{\partial y}\frac{\partial y}{\partial r}$.

Similar logic for $s$ and $t$.
Looking at the options in the PDF (Screenshot 9), the correct set is the second option circle:
\[ \frac{\partial u}{\partial r} = \frac{\partial u}{\partial x}\frac{\partial x}{\partial r} + \frac{\partial u}{\partial y}\frac{\partial y}{\partial r}, \dots \]
(Note: The first option has $\partial u/\partial r$ on the RHS which is wrong. The third option adds a third term. The correct option matches the standard 2-intermediate variable form).

\textbf{Final Answer:} Select the option:
$\frac{\partial u}{\partial r} = \frac{\partial u}{\partial x}\frac{\partial x}{\partial r} + \frac{\partial u}{\partial y}\frac{\partial y}{\partial r}$, etc. (Usually the 2nd bubble).

% -------------------------------------------------------

\subsection*{Problem 10: Evaluation at a Point}
\textbf{Given:} $z = x^4 + x^2y$, $x = s + 2t - u$, $y = stu^2$.
Point: $s=1, t=2, u=4$.

Step 1: Calculate $x$ and $y$ at this point.
$x = 1 + 2(2) - 4 = 1 + 4 - 4 = 1$.
$y = (1)(2)(4^2) = 2(16) = 32$.

Step 2: Calculate Partials of $z$ at $(x,y)=(1,32)$.
$\frac{\partial z}{\partial x} = 4x^3 + 2xy = 4(1)^3 + 2(1)(32) = 4 + 64 = 68$.
$\frac{\partial z}{\partial y} = x^2 = 1^2 = 1$.

Step 3: Calculate Partials of $x, y$ w.r.t $s, t, u$.
$x_s=1, x_t=2, x_u=-1$.
$y_s=tu^2, y_t=su^2, y_u=2stu$.

Evaluate these at $(1,2,4)$:
$y_s = (2)(16) = 32$.
$y_t = (1)(16) = 16$.
$y_u = 2(1)(2)(4) = 16$.

Step 4: Chain Rule Assembly.
$\frac{\partial z}{\partial s} = z_x x_s + z_y y_s = (68)(1) + (1)(32) = 100$.
$\frac{\partial z}{\partial t} = z_x x_t + z_y y_t = (68)(2) + (1)(16) = 136 + 16 = 152$.
$\frac{\partial z}{\partial u} = z_x x_u + z_y y_u = (68)(-1) + (1)(16) = -68 + 16 = -52$.

\textbf{Final Answers:}
$\frac{\partial z}{\partial s} = \boxed{100}, \quad \frac{\partial z}{\partial t} = \boxed{152}, \quad \frac{\partial z}{\partial u} = \boxed{-52}$

% -------------------------------------------------------

\subsection*{Problem 11: Repetition with New Values}
\textbf{Given:} Same function as \#10.
Point: $s=4, t=2, u=5$.

Step 1: Find $x, y$.
$x = 4 + 2(2) - 5 = 4 + 4 - 5 = 3$.
$y = (4)(2)(5^2) = 8(25) = 200$.

Step 2: Find $z_x, z_y$.
$z_x = 4x^3 + 2xy = 4(27) + 2(3)(200) = 108 + 1200 = 1308$.
$z_y = x^2 = 9$.

Step 3: Find deriv values.
$y_s = tu^2 = 2(25) = 50$.
$y_t = su^2 = 4(25) = 100$.
$y_u = 2stu = 2(4)(2)(5) = 80$.
$x_s=1, x_t=2, x_u=-1$.

Step 4: Assemble.
$\frac{\partial z}{\partial s} = (1308)(1) + (9)(50) = 1308 + 450 = 1758$.
$\frac{\partial z}{\partial t} = (1308)(2) + (9)(100) = 2616 + 900 = 3516$.
$\frac{\partial z}{\partial u} = (1308)(-1) + (9)(80) = -1308 + 720 = -588$.

\textbf{Final Answers:} $\boxed{1758}, \boxed{3516}, \boxed{-588}$

% -------------------------------------------------------

\subsection*{Problem 12: Implicit Differentiation}
\textbf{Equation:} $\cos(xy) = 3 + \sin(y)$.
Rewrite as $F(x,y) = 0$:
$F(x,y) = \cos(xy) - \sin(y) - 3 = 0$.

Formula: $\frac{dy}{dx} = -\frac{F_x}{F_y}$.

$F_x = \frac{\partial}{\partial x}[\cos(xy) - \sin(y) - 3] = -\sin(xy) \cdot y = -y\sin(xy)$.
$F_y = \frac{\partial}{\partial y}[\cos(xy) - \sin(y) - 3] = -\sin(xy) \cdot x - \cos(y) = -x\sin(xy) - \cos(y)$.

Apply formula:
\[ \frac{dy}{dx} = -\frac{-y\sin(xy)}{-x\sin(xy) - \cos(y)} = \frac{y\sin(xy)}{-(x\sin(xy) + \cos(y))} = -\frac{y\sin(xy)}{x\sin(xy) + \cos(y)} \]

\textbf{Final Answer:} $\boxed{-\frac{y\sin(xy)}{x\sin(xy) + \cos(y)}}$

% -------------------------------------------------------

\subsection*{Problem 13: Implicit Diff for Surfaces}
\textbf{Equation:} $x^2 - y^2 + z^2 - 8z = 9$.
$F(x,y,z) = x^2 - y^2 + z^2 - 8z - 9 = 0$.

Formulas: $\frac{\partial z}{\partial x} = -\frac{F_x}{F_z}, \quad \frac{\partial z}{\partial y} = -\frac{F_y}{F_z}$.

$F_x = 2x$.
$F_y = -2y$.
$F_z = 2z - 8$.

$\frac{\partial z}{\partial x} = -\frac{2x}{2z - 8} = -\frac{2x}{2(z-4)} = -\frac{x}{z-4} = \frac{x}{4-z}$.
$\frac{\partial z}{\partial y} = -\frac{-2y}{2z - 8} = \frac{2y}{2(z-4)} = \frac{y}{z-4}$.

\textbf{Final Answers:}
$\frac{\partial z}{\partial x} = \boxed{\frac{x}{4-z}}$
$\frac{\partial z}{\partial y} = \boxed{\frac{y}{z-4}}$

% -------------------------------------------------------

\subsection*{Problem 14: Rate of Change (Word Problem)}
\textbf{Given:} $T(x,y)$. $T_x(5,8)=8, T_y(5,8)=6$.
Paths: $x(t) = \sqrt{3+t}$, $y(t) = 7 + \frac{1}{22}t$.
Find $dT/dt$ at $t=22$.

Step 1: Check location at $t=22$.
$x(22) = \sqrt{3+22} = \sqrt{25} = 5$.
$y(22) = 7 + 22/22 = 7+1 = 8$.
Location is $(5,8)$, which matches the given partials.

Step 2: Differentiate paths.
$x'(t) = \frac{1}{2\sqrt{3+t}}$. At $t=22$: $x'(22) = \frac{1}{2(5)} = \frac{1}{10} = 0.1$.
$y'(t) = \frac{1}{22}$.

Step 3: Chain Rule.
\[ \frac{dT}{dt} = T_x x'(t) + T_y y'(t) \]
\[ \frac{dT}{dt} = (8)\left(\frac{1}{10}\right) + (6)\left(\frac{1}{22}\right) \]
\[ \frac{dT}{dt} = \frac{8}{10} + \frac{6}{22} = \frac{4}{5} + \frac{3}{11} \]
Common denominator 55:
\[ \frac{44}{55} + \frac{15}{55} = \frac{59}{55} \]

\textbf{Final Answer:} $\boxed{59/55}$

% -------------------------------------------------------

\subsection*{Problem 15 & 16: Second Derivatives (The Challenge Problem)}
\textbf{Given:} $z = f(x,y)$, $x = 2r\cos\theta$, $y = 2r\sin\theta$.

\textbf{Part (a): Find $\frac{\partial z}{\partial r}$}
$x_r = 2\cos\theta$, $y_r = 2\sin\theta$.
\[ \frac{\partial z}{\partial r} = \frac{\partial z}{\partial x}(2\cos\theta) + \frac{\partial z}{\partial y}(2\sin\theta) \]
Box format from screenshot:
$\frac{\partial z}{\partial r} = \frac{\partial z}{\partial x} \boxed{2\cos\theta} + \frac{\partial z}{\partial y} \boxed{2\sin\theta}$.

\textbf{Part (b): Find $\frac{\partial z}{\partial \theta}$}
$x_\theta = -2r\sin\theta$, $y_\theta = 2r\cos\theta$.
$\frac{\partial z}{\partial \theta} = \frac{\partial z}{\partial x} \boxed{-2r\sin\theta} + \frac{\partial z}{\partial y} \boxed{2r\cos\theta}$.

\textbf{Part (c): Find $\frac{\partial^2 z}{\partial r \partial \theta}$}
We differentiate the result of (a) with respect to $\theta$.
\[ \frac{\partial}{\partial \theta} \left[ \frac{\partial z}{\partial x}(2\cos\theta) + \frac{\partial z}{\partial y}(2\sin\theta) \right] \]
Apply Product Rule to each term:
Term 1: $\frac{\partial}{\partial \theta}(z_x \cdot 2\cos\theta) = z_x(-2\sin\theta) + (2\cos\theta)\frac{\partial}{\partial \theta}(z_x)$.
Term 2: $\frac{\partial}{\partial \theta}(z_y \cdot 2\sin\theta) = z_y(2\cos\theta) + (2\sin\theta)\frac{\partial}{\partial \theta}(z_y)$.

Now expand $\frac{\partial}{\partial \theta}(z_x)$ using Chain Rule:
$\frac{\partial z_x}{\partial \theta} = \frac{\partial z_x}{\partial x}\frac{\partial x}{\partial \theta} + \frac{\partial z_x}{\partial y}\frac{\partial y}{\partial \theta} = z_{xx}(-2r\sin\theta) + z_{xy}(2r\cos\theta)$.
Similarly for $z_y$:
$\frac{\partial z_y}{\partial \theta} = z_{yx}(-2r\sin\theta) + z_{yy}(2r\cos\theta)$.

Combine everything:
Result $= -2\sin\theta z_x + 2\cos\theta [ z_{xx}(-2r\sin\theta) + z_{xy}(2r\cos\theta) ] + 2\cos\theta z_y + 2\sin\theta [ z_{yx}(-2r\sin\theta) + z_{yy}(2r\cos\theta) ]$.

Group by partials (assuming $z_{xy} = z_{yx}$):
$z_x$ coeff: $-2\sin\theta$
$z_y$ coeff: $2\cos\theta$
$z_{xx}$ coeff: $2\cos\theta(-2r\sin\theta) = -4r\sin\theta\cos\theta$.
$z_{yy}$ coeff: $2\sin\theta(2r\cos\theta) = 4r\sin\theta\cos\theta$.
$z_{xy}$ coeff: $2\cos\theta(2r\cos\theta) + 2\sin\theta(-2r\sin\theta) = 4r(\cos^2\theta - \sin^2\theta)$.

The screenshot asks to fill in boxes:
$= \frac{\partial z}{\partial y} \boxed{2\cos\theta} + \boxed{-2\sin\theta} \frac{\partial z}{\partial x} + \boxed{4r\cos^2\theta - 4r\sin^2\theta} \frac{\partial^2 z}{\partial y \partial x} + \boxed{4r\sin\theta\cos\theta} \left( \frac{\partial^2 z}{\partial y^2} - \frac{\partial^2 z}{\partial x^2} \right)$.

Wait, let's verify the grouping for the last term to match the screenshot form:
The screenshot has: $\dots + (\text{Box}) \left( \frac{\partial^2 z}{\partial y^2} - \frac{\partial^2 z}{\partial x^2} \right)$.
My $z_{yy}$ term was $4r\sin\theta\cos\theta$. My $z_{xx}$ term was $-4r\sin\theta\cos\theta$.
Factor out $4r\sin\theta\cos\theta$:
$4r\sin\theta\cos\theta (z_{yy} - z_{xx})$. This matches perfectly.

\textbf{Final Answers for Filling Blanks:}
First box ($z_y$ coeff): $\boxed{2\cos\theta}$
Second box ($z_x$ coeff): $\boxed{-2\sin\theta}$
Third box (Mixed partial coeff): $\boxed{4r(\cos^2\theta - \sin^2\theta)}$ or $\boxed{4r\cos(2\theta)}$
Fourth box (Difference coeff): $\boxed{4r\sin\theta\cos\theta}$ or $\boxed{2r\sin(2\theta)}$

\section{Part 3: In-Depth Analysis}

\subsection{Problem Types}
\begin{enumerate}
    \item \textbf{Type 1: Direct Computation (Problems 1, 3, 4, 6).} Given specific functions, find the derivative w.r.t $t$. Strategy: Compute partials, compute path derivatives, multiply and sum.
    \item \textbf{Type 2: Multivariable Computation (Problems 2, 5, 7, 10, 11).} The path depends on two variables ($s,t$). Strategy: Treat $s$ as constant when finding $\partial/\partial t$, and vice-versa.
    \item \textbf{Type 3: Implicit Differentiation (Problems 12, 13).} $F(x,y)=0$ or $F(x,y,z)=0$. Strategy: Move all terms to one side, define $F$, use ratio of partials formula.
    \item \textbf{Type 4: Data Table/Abstract (Problem 8, 14).} No functions given, only values at specific points. Strategy: Set up the Chain Rule formula symbolically first, then plug in numbers.
    \item \textbf{Type 5: Second Derivative Transformation (Problem 15).} Coordinate transformation (Polar). Strategy: Use Product Rule on the first derivative result, applying Chain Rule again to the inner partials.
\end{enumerate}

\subsection{Key Techniques}
\begin{itemize}
    \item \textbf{The "Evaluate Last" Trick:} In Problems 10 and 11, it is much faster to find the numerical values of $x, y$ and the derivatives first, rather than substituting the algebraic expressions of $s,t,u$ into $z$. Algebra is prone to error; arithmetic is safer.
    \item \textbf{The Inverse Tangent Derivative:} In Problem 4, knowing $\frac{d}{dt}\arctan(t) = \frac{1}{1+t^2}$ was crucial.
    \item \textbf{Implicit Formula Shortcut:} Instead of differentiating term-by-term implicitly (Calc I style), defining $F(x,y,z)$ and using $-F_x/F_z$ (Problem 13) saves massive amounts of algebraic rearrangement.
\end{itemize}

\section{Part 4: Cheatsheet}

\subsection{Formulas}
\begin{itemize}
    \item \textbf{Chain Rule:} $\frac{dz}{dt} = \nabla z \cdot \mathbf{r}'(t) = \langle z_x, z_y \rangle \cdot \langle x', y' \rangle$.
    \item \textbf{Implicit (2D):} $\frac{dy}{dx} = -\frac{F_x}{F_y}$.
    \item \textbf{Implicit (3D):} $\frac{\partial z}{\partial x} = -\frac{F_x}{F_z}, \quad \frac{\partial z}{\partial y} = -\frac{F_y}{F_z}$.
\end{itemize}

\subsection{Common Pitfalls}
\begin{enumerate}
    \item \textbf{Evaluating at $t$ instead of $x$:} In Problem 14, you are given $T_x(5,8)$. A common error is plugging $t=22$ into $T_x$ directly, but $T_x$ takes spatial coordinates, not time.
    \item \textbf{Sign Errors in Implicit Diff:} The formula has a negative sign: $-F_x/F_y$. Students often forget it.
    \item \textbf{Mixing Variables:} In the final answer, if the question asks for "in terms of $t$," you must replace every $x$ and $y$. If it asks for "in terms of $x$ and $y$," leave them.
\end{enumerate}

\section{Part 5: Conceptual Synthesis}

\subsection{Thematic Connections}
The Chain Rule represents the \textbf{Linearization of Composition}. Just as a tangent plane approximates a surface, the Chain Rule tells us that the total change is the linear combination of changes in each direction. It connects calculus to linear algebra: the derivative is a linear transformation (matrix multiplication of the Jacobian).

\subsection{Forward and Backward Links}
\begin{itemize}
    \item \textbf{Backward:} Relies heavily on Calc I Chain Rule and Calc III Partial Derivatives.
    \item \textbf{Forward:} Essential for \textbf{Directional Derivatives} (next chapter), \textbf{Gradient Vectors}, and \textbf{Change of Variables in Multiple Integrals} (Jacobians). You cannot do Triple Integrals in Spherical Coordinates without understanding this logic.
\end{itemize}

\section{Part 6: Real-World Application (Finance)}

\subsection{Scenario: Derivatives Pricing (The "Greeks")}
In Quantitative Finance, the Black-Scholes equation models the price of an option $V(S, t, \sigma, r)$, where $S$ is stock price, $t$ is time, $\sigma$ is volatility, and $r$ is interest rate.
Traders need to know how the value of their portfolio changes over time. However, $S$ changes with time (stochastically), and $\sigma$ might change with time.
To find $\frac{dV}{dt}$ (total time decay and drift), they use the Multivariable Chain Rule (Itô's Lemma is the stochastic calculus version of this).
\begin{itemize}
    \item $\frac{\partial V}{\partial S}$ is called \textbf{Delta} ($\Delta$).
    \item $\frac{\partial V}{\partial t}$ is called \textbf{Theta} ($\Theta$).
    \item The Chain Rule allows traders to "hedge" risks by canceling out specific partial derivatives.
\end{itemize}

\subsection{Model Setup}
Let $P(S, r)$ be the value of a bond portfolio.
$S(t) = 100 + e^{0.05t}$ (Stock index growing).
$r(t) = 0.02 + 0.001\sin(t)$ (Interest rates fluctuating).
Find the rate of change of the portfolio value $\frac{dP}{dt}$.
\[ \frac{dP}{dt} = \frac{\partial P}{\partial S}\frac{dS}{dt} + \frac{\partial P}{\partial r}\frac{dr}{dt} \]
This determines if the portfolio is making or losing money at an exact instant.

\section{Part 7: Common Variations and Untested Concepts}

\subsection{Homogeneous Functions (Euler's Theorem)}
\textit{Not covered in homework.}
If $f(x,y)$ is homogeneous of degree $n$ (i.e., $f(tx, ty) = t^n f(x,y)$), then:
\[ x\frac{\partial f}{\partial x} + y\frac{\partial f}{\partial y} = n f(x,y) \]
\textbf{Example:} $f(x,y) = x^2 + y^2$ (degree 2).
$f_x = 2x, f_y = 2y$.
$x(2x) + y(2y) = 2x^2 + 2y^2 = 2(x^2+y^2) = 2f$. Verified.

\section{Part 8: Advanced Diagnostic Testing ("Find the Flaw")}

\subsection*{Flaw 1: The Sign Trap}
\textbf{Problem:} Find $\frac{dy}{dx}$ if $x^2 + y^2 - 2y = 0$.
\textbf{Flawed Solution:}
$F_x = 2x$. $F_y = 2y - 2$.
$\frac{dy}{dx} = \frac{F_x}{F_y} = \frac{2x}{2y-2} = \frac{x}{y-1}$.
\textbf{Analysis:} The student forgot the negative sign in the implicit differentiation formula.
\textbf{Correction:} $\frac{dy}{dx} = -\frac{F_x}{F_y} = -\frac{x}{y-1}$.

\subsection*{Flaw 2: The Variable Mix-Up}
\textbf{Problem:} $z = x^2y$, $x=t^2, y=t^3$. Find $dz/dt$ at $t=1$.
\textbf{Flawed Solution:}
$z_x = 2xy$. At $t=1$, $x=1, y=1 \to z_x = 2$.
$z_y = x^2$. At $t=1$, $z_y = 1$.
$dz/dt = z_x + z_y = 2 + 1 = 3$.
\textbf{Analysis:} The student summed the partials but forgot to multiply by the derivatives of the paths ($dx/dt, dy/dt$).
\textbf{Correction:} $x'=2t=2, y'=3t^2=3$. $dz/dt = 2(2) + 1(3) = 7$.

\subsection*{Flaw 3: Product Rule Negligence}
\textbf{Problem:} $\partial^2/\partial t \partial x$ if $z = f(x, t)$ and $x = g(t)$.
\textbf{Flawed Solution:} $\frac{\partial}{\partial t}(\frac{\partial z}{\partial x}) = \frac{\partial^2 z}{\partial t \partial x}$.
\textbf{Analysis:} Since $z$ depends on $x$ and $x$ depends on $t$, taking $\partial/\partial t$ requires the Chain Rule again.
\textbf{Correction:} $\frac{\partial}{\partial t}(z_x) = z_{xt} + z_{xx}\frac{dx}{dt}$.

\subsection*{Flaw 4: Incorrect Tree Diagram}
\textbf{Problem:} $w = f(x,y,z)$, $x=g(u,v)$, $y=h(u,v)$. Find $\partial w / \partial u$.
\textbf{Flawed Solution:} $\frac{\partial w}{\partial u} = \frac{\partial w}{\partial x}\frac{\partial x}{\partial u} + \frac{\partial w}{\partial y}\frac{\partial y}{\partial u} + \frac{\partial w}{\partial z}$.
\textbf{Analysis:} The term $\frac{\partial w}{\partial z}$ is added as a "loose end." If $z$ does not depend on $u$, its contribution is zero. If $z$ is independent, that branch doesn't exist.
\textbf{Correction:} Remove the last term.

\subsection*{Flaw 5: Implicit Diff Algebra}
\textbf{Problem:} Find $y'$ for $xy = e^y$.
\textbf{Flawed Solution:} Diff both sides: $y + xy' = e^y$. Solve: $xy' = e^y - y \to y' = \frac{e^y-y}{x}$.
\textbf{Analysis:} The RHS derivative of $e^y$ is $e^y y'$, not just $e^y$.
\textbf{Correction:} $y + xy' = e^y y' \Rightarrow y = y'(e^y - x) \Rightarrow y' = \frac{y}{e^y - x}$.

\end{document}